\chapter{Contexte general du projet}
\label{chap:contexte}

\section{Organisme d'accueil}

Ce projet a ete realise dans le cadre d'un stage de fin d'annee au sein de
\organisme\footnote{\acompleter{A completer : presentation de l'agence/departement
d'accueil, activites, effectif, positionnement sur le marche de la cybersecurite et
des services IT.}}. \organisme{} est une entreprise de services numeriques qui
accompagne ses clients dans la modernisation, la securisation et l'exploitation de
leurs systemes d'information. Le stage s'est deroule au sein d'une equipe technique
en charge de la securite operationnelle (\emph{DevSecOps}), dont l'une des missions
est de fournir aux equipes clientes une visibilite fiable et exploitable sur
l'exposition aux vulnerabilites de leur parc informatique.

\section{Problematique}
\label{sec:problematique}

Le point de depart du projet est un constat operationnel, formalise dans les notes de
cadrage initiales (voir \emph{app\_notes\_dxc.txt}) : les donnees necessaires a
l'evaluation du risque existent, mais elles sont fragmentees et de natures tres
differentes.

\begin{itemize}
  \item Les exports \textbf{Qualys} decrivent, pour chaque hote scanne, la liste des
  vulnerabilites detectees (identifiant QID, severite, score QDS, dates de premiere
  et derniere detection, statut de correction). Ces donnees sont \emph{volumineuses}
  et \emph{mises a jour frequemment} (un nouveau scan peut redefinir l'etat complet
  du parc).
  \item La \textbf{CMDB} (Configuration Management Database) decrit, pour chaque
  element de configuration, son role technique, son environnement (production, test,
  developpement...), son proprietaire metier et les mesures de protection deployees
  (antivirus/EDR, supervision SOC, cycle de patch standard). Ces donnees sont
  \emph{plus stables} mais \emph{lues frequemment} par tous les traitements de
  scoring.
\end{itemize}

Avant ce projet, le rapprochement entre ces deux sources -- indispensable pour
distinguer une vulnerabilite critique sur un serveur de production expose d'une
vulnerabilite identique sur un poste de test isole -- se faisait manuellement (feuille
de calcul), ce qui est lent, source d'erreurs, non tracable dans le temps et ne
permet aucune restitution visuelle consolidee pour le pilotage. Cette absence de
score de risque contextualise et de gouvernance SLA associee constitue la
problematique centrale adressee par ce projet.

\section{Objectifs du projet}

Les objectifs assignes au stage, tels qu'ils decoulent du cadrage initial, sont les
suivants :

\begin{enumerate}
  \item Concevoir un modele mathematique de risque qui combine la criticite metier
  d'un actif (issue de la CMDB), la gravite intrinseque de ses vulnerabilites (issue
  de Qualys), son exposition reseau et ses controles compensatoires, en un score
  unique, explicable et reproductible : l'\emph{Asset Risk Score} (ARS).
  \item Automatiser l'ingestion des exports Qualys et CMDB (fichiers XML), avec
  detection des doublons et traitement asynchrone, pour supporter des mises a jour
  frequentes sans intervention manuelle.
  \item Deriver de l'ARS une gouvernance operationnelle : classification en 5 classes
  de risque (\emph{Green} a \emph{Black}), delais de remediation SLA associes, et
  gestion des exceptions de risque accepte.
  \item Mesurer objectivement la performance des equipes de remediation via un
  indicateur de temps moyen de correction pondere par criticite (MTTR pondere).
  \item Restituer l'ensemble de ces informations dans un tableau de bord temps reel,
  exploitable par un RSSI ou un analyste securite, combinant plusieurs types de
  graphiques (repartition, flux, historique, conformite SLA).
  \item Offrir un canal d'interrogation en langage naturel de ces donnees via un
  assistant conversationnel fonde sur une architecture RAG.
\end{enumerate}

\section{Perimetre du projet}
\label{sec:perimetre}

\subsection{Ce qui est couvert}

\begin{itemize}
  \item L'ingestion, le parsing et la persistance des exports Qualys (rapport de
  vulnerabilites par hote) et CMDB (inventaire des elements de configuration), au
  format XML.
  \item Le calcul, pour chaque actif reconcilie (present a la fois dans Qualys et
  dans la CMDB), de l'ACS, du score QDS pondere par severite, de l'ARS, de la classe
  de risque, du statut SLA et de l'anciennete de la vulnerabilite ouverte la plus
  ancienne.
  \item La configuration, par l'utilisateur, des poids du modele : poids par
  severite, poids par role technique et par environnement (pour l'ACS), facteurs de
  modulation ($I_{External}$, $M_{SOC}$, $M_{AV}$, $M_{Patch}$), seuils et delais SLA
  par classe, poids de ponderation du MTTR par classe.
  \item Un tableau de bord de pilotage alimente en temps reel (Server-Sent Events)
  presentant la repartition du parc par classe de risque, par equipe proprietaire,
  les flux Environnement/Role/Classe (diagramme de Sankey), l'historique de l'ARS
  moyen et la conformite SLA (age des vulnerabilites ouvertes vs. ARS).
  \item Un assistant conversationnel RAG capable de repondre, en langage naturel, a
  des questions portant sur les donnees ingerees et calculees -- livre sous forme de
  prototype fonctionnel, ce module restant au stade du developpement et de la
  conception (cf. limites, \cref{sec:perspectives}).
\end{itemize}

\subsection{Ce qui est explicitement hors perimetre}

\begin{itemize}
  \item L'agregation multi-scan dans le temps au niveau de la CMDB : chaque nouvel
  import remplace integralement l'etat precedent (strategie \emph{wipe-and-reingest},
  cf. \cref{sec:pipeline-ingestion}) ; il n'y a pas, dans la version realisee, de
  gestion de versions historiques des CI.
  \item L'authentification et le controle d'acces par role (RBAC) : la version
  realisee est un prototype fonctionnel dont la chaine de securite HTTP est
  volontairement ouverte (voir limites, \cref{sec:perspectives}).
  \item La connexion a d'autres sources de vulnerabilites que Qualys (Nessus,
  Tenable, OpenVAS...) ou a d'autres CMDB que le format d'export etudie.
  \item Le calcul effectif du facteur de penalite operationnelle
  $\mathrm{Global\_I_{Ops}}$ dans l'agregation \emph{Enterprise Risk Score} : le
  modele mathematique est entierement specifie et demontre (\cref{chap:modelisation})
  mais son integration dans le service applicatif reste une perspective
  d'evolution.
\end{itemize}

\section{Methodologie de travail}

Le developpement a suivi une demarche iterative et incrementale, guidee par
l'historique des livraisons du projet : une premiere version couvrant l'ingestion et
le modele de donnees de base, suivie d'une phase de fiabilisation de l'ingestion
CMDB, puis d'une extension fonctionnelle majeure (moteur de risque configurable,
tableau de bord temps reel, assistant RAG), et enfin d'une phase de
conteneurisation en vue du deploiement (Docker Compose). Chaque etape a fait l'objet
de commits Git dedies, permettant une tracabilite complete des choix
d'implementation. Cette approche a permis de valider tot le pipeline de bout en
bout (upload $\to$ ingestion $\to$ calcul $\to$ restitution) avant d'en enrichir
progressivement la richesse fonctionnelle (configuration dynamique des poids,
gestion des exceptions, historisation des instantanes de risque).

\section{Organisation du reste du rapport}

La suite du rapport suit la logique amont-aval d'un projet logiciel : specification
des besoins (\cref{chap:analyse}), conception (\cref{chap:conception}), fondements
mathematiques du coeur metier (\cref{chap:modelisation}), realisation
(\cref{chap:realisation}), restitution visuelle (\cref{chap:visualisation}) et
validation (\cref{chap:tests}), avant de conclure sur le bilan et les perspectives.
